\section{结论}
\label{sec:conclusion}

我们对 \pol{} 的结论是明确的，并且依现有证据，对每一项作出决定的功能均为否定（\cref{sec:verdict}）。
托管的 \jev{} 及与其一同测试的类型化模型，以及出于审慎原则，与它们一同测试的低成本生成式评估器（尚未表明它们能避免这些失效，且其中多数失效从未经过比较），都不应依 \art{10} 裁决争议、依 \art{6}(3) 自动进行伦理验证、依 \clause{5.6} 量化个人，或使对某人言论的阻断成为最终决定；原因在于：阈值无法迁移，判定可能被其所审计的内容所操纵，并可能取决于选项名称而非定义，明确的“不知道”选项在其为正确答案时很少被选择，而“以上皆非”的选择不一致，并且此类模型所要升级至的同类评估器几乎重复了其全部高置信错误；对于逐人量化，这一结论还基于审慎原则，因为尚无研究按说话者群体测量误差。
对于检测类功能，即 \clause{5.4.1} 的安全阀、\clause{4.3.3.1} 的截断以及 \art{7} 针对文本的异常检测（若无经过训练的模型，则不适用于数值或结构化信号），托管的 \jev{} 只能在\cref{sec:design}的设计中作为有记录、三值、经本地校准的\emph{检测器}使用；在该设计中，自动行动须以独立构建的检测器之间取得一致为前提，对个人的任何限制均为短暂、可逆的，每一条可能导致阻断的链路都以人为终点，高置信的自动决策须受审计，且每一项决策均予记录；\clause{5.6} 的公共决策评估也只能在同样意义上使用，即作为一种经审计的汇总测量，并将其测得误差带入每一项结论。
即便是这一检测器结论，也是从二值仇恨言论的英文基准外推而来的：尚无研究在中文中测量该规范的三值构念，所找到的唯一一项中文情绪结果显示所有系统均表现较弱，而所测试的每一个系统对仇恨言论与仅具冒犯性的语言的区分能力都较弱。
支撑这些结论的每一项发现的确定性均为低或极低，其中四项发现有争议，截至 10 月 8 日的更新检索未推翻其中任何一项；一项针对该规范自身构念、设计良好的研究即可改变其中任何一项。

尽管如此，类型化决策模型的出现带来了自动化治理保障机制所需的接口：对每一次输入与输出给出带概率的类型化判定，且成本低到足以支撑全面筛查。
无论采用何种模型，本评估中可借鉴到 \pol{} 的，是一组与该接口相关的做法：三值问题，并单独询问充分性；绑定到定义并经重命名检验的中性标识符；记录每一项决策；本地校准并进行时间再验证；以构造不同的检测器之间取得一致作为门控；以人为终点的链路，以及对高置信决策的盲法审计；不向不受信任的调用方提供概率；不需要的敏感字段不放入输入，不受信任的查询受到限制并予记录；独立标记在不同调用中分别询问；对事件而非对人评分；以及，在上述任何一项能够被测量而非外推之前，一个中文构念基准。

就 \pol{} 而言，该框架自身的若干承诺与证据所建议的做法一致：不在场状态、关于爱与恨并非道德标签的告诫，以及对人与行为的区分。
另一些条款则被证据凸显为文本尚待作出的选择：是否在理解语言中的情绪与推断个人的情绪状态之间划出明确界线；如何向不能可靠地回答“信息不足”的模型提出三值问题；如何协调透明性与封闭模型；对于具有实质后果的决定，非干预原则是否应让位于能够暂停这些决定的人，这接近于数据保护法在其适用范围内的要求；以及逐人评估能否与平等联结相调和。
将这些选择明确化，并在该规范自身的构念上检验由此形成的保障机制，是后续的工作。
